\section{The Subgraph With No Service}
\label{sec:ch14-the-subgraph-with-no-service}
\index{event-driven federated subscriptions|(}%

Everything so far has taken one arrangement for granted: a subscription is a
field of a service, and the router holds a connection to that service for as
long as a client is listening. Cosmo offers a second arrangement in which none
of that is true, and the fastest way to say what it is is to show what Mosaic
had to add.

\begin{graphqlsdl}[fontsize=\footnotesize]
directive @edfs__natsSubscribe(
  subjects: [String!]!
  providerId: String! = "default"
  streamConfiguration: edfs__NatsStreamConfiguration
) on FIELD_DEFINITION

type Subscription {
  reviewPublished(productId: ID!): Review!
    @edfs__natsSubscribe(subjects: ["mosaic.reviewAdded.{{ args.productId }}"])
}

type Review @key(fields: "id", resolvable: false) {
  id: ID! @external
}
\end{graphqlsdl}

That is most of \code{schema/streams.graphql}. What the listing leaves out is
the \code{@link} block every subgraph in this book carries and the
\code{edfs\_\_NatsStreamConfiguration} input the directive's third argument
refers to, which brings the file to twenty-eight lines. There is no project
beside it, no port, no process. The documentation is direct about this: an
event-driven subgraph \enquote{does not need to be implemented. It is simply a
Subgraph Schema that instructs the Router on how to connect specific root
fields to the Event Source}~\autocite{cosmodocs2026streams}.

Note that the directive is declared here rather than assumed. The composer has
a built-in definition for it, as it does for \code{@key}, but a subgraph schema
is published SDL and has to stand on its own, so the file spells out the
directive it uses.

Chapter~\ref{ch:hard-modeling-problems} gave Mosaic a service that owns no
domain. This is one better: a subgraph that owns no process.

\subsection*{What the composer makes of it}

\index{composition!a datasource that is not GraphQL}%
Add it to \code{federation/mosaic.yaml} as an eighth entry and compose. The
routing table gains a datasource, and it is not the kind every other entry in
this book has been:

\begin{minted}[fontsize=\scriptsize,breaklines]{json}
{
  "kind": "PUBSUB",
  "customEvents": {"nats": [{
    "engineEventConfiguration": {"providerId": "default", "type": "SUBSCRIBE",
                                 "typeName": "Subscription", "fieldName": "reviewPublished"},
    "subjects": ["mosaic.reviewAdded.{{ args.productId }}"]}]}
}
\end{minted}

Seven \code{GRAPHQL} sources with a fetch url each, and one \code{PUBSUB}
source with a broker subject and no url anywhere in it. The entry in the input
file was still obliged to give a \code{routing\_url}, because the schema
requires one, and the composer does copy it into the config's subgraph list
beside the other seven. What it does not do is give this datasource anything to
fetch with, so nothing ever calls it. Mosaic points it at 5108, where nothing
listens, so that a reader who suspects something is quietly serving this can
check.

The template in the subject is the other half. \code{\{\{ args.productId \}\}}
is interpolated from the subscription document, so a client subscribing to one
product makes the router listen on one subject.

\subsection*{Where the broker lives}

\index{router!the \code{events} block}%
The schema says which subject; it does not say where the broker is. That is a
separate block in \code{router/config.yaml}:

\begin{minted}[fontsize=\footnotesize,breaklines]{yaml}
events:
  providers:
    nats:
      - id: default
        url: "nats://host.docker.internal:4222"
\end{minted}

and the split is deliberate on Cosmo's part rather than an accident of
layout. A subject name is part of a contract between two teams and belongs
where a composer can read it. A broker url and its credentials are a property
of one deployment and belong with the rest of that deployment's addresses.
\code{default} is the provider id because \code{providerId} on the directive
defaults to that string and nothing overrides it; a second broker would be a
second entry here and an argument on the directive.

On startup the router says so:

\begin{minted}[fontsize=\footnotesize,breaklines]{text}
INFO core/router.go:546 Nats Event source enabled {"provider_id": "default"}
INFO nats/provider_builder.go:87 NATS connection established
     {"provider_id": "default", "url": "nats://host.docker.internal:4222"}
\end{minted}

\subsection*{What goes on the wire is almost nothing}

\index{event-driven federated subscriptions!the message is a representation}%
Here is the part that makes the feature worth the trouble. Publish this to
NATS, by hand, with any client:

\begin{minted}[fontsize=\footnotesize,breaklines]{text}
subject  mosaic.reviewAdded.UHJvZHVjdDoAAACgAAAAQIAAAAAAAAAW
body     {"__typename":"Review","id":"UmV2aWV3OgAAAOAAAABAgAAAAAAAAIc="}
\end{minted}

and this arrives at a client subscribed through the router:

\begin{minted}[fontsize=\scriptsize,breaklines]{text}
event: next
data: {"data":{"reviewPublished":{"id":"UmV2aWV3OgAAAOAAAABAgAAAAAAAAIc=",
       "rating":5,"body":"End grain, so it is kind to the knife edge.",
       "author":{"displayName":"Katharina Brandt"}}}}
\end{minted}

A type name and a key went in. Fields from two services came out, neither of
them the one that declared the field. The router treats the message body as a
federation representation, exactly like the ones
chapter~\ref{ch:entity-resolution-done-right} watched being built from a query
plan, and resolves the rest of the payload the ordinary way. The plan for that
document says as much:

\begin{minted}[fontsize=\scriptsize,breaklines]{text}
Sequence
  Trigger  7-nats-0  fetchId 0
  Entity   reviews   fetchId 1  dependsOnFetchIds [0]
           ... on Review { __typename rating body author { __typename id } }
  Entity   accounts  fetchId 2  dependsOnFetchIds [1]
           ... on Customer { __typename displayName }
\end{minted}

Two entity fetches where the held subscription had one, because the message
carries no data at all: even \code{rating}, which Reviews owns and which the
publisher had in its hand, has to be fetched back. And the trigger names
\code{7-nats-0}, which is a datasource position and a provider rather than a
subgraph, because there is no service to name. That string is a positional
artefact rather than a stable identifier, which is why the Postman assertion in
this chapter's collection checks only that the trigger's source mentions the
broker.

\index{event-driven federated subscriptions!why the payload is not published}%
The obvious objection is that publishing the whole review would save a hop, and
the obvious answer is that it would publish a second copy of a schema Reviews
owns. Every field added to \code{Review} would then have to be added to the
message, every consumer would be reading a shape nothing validates, and the
first field somebody forgot would go stale silently. What is on the wire is a
pointer, and a pointer cannot drift.

Figure~\ref{fig:ch14-two-arrangements} puts the two arrangements beside each
other, which is the comparison the whole chapter has been building towards.

\begin{figure}[htbp]
  \centering
  % The two subscriptions, from one write. Referenced from chapter 14. Bare
% tikzpicture; the figure environment, caption and label live at the call site.
%
% The point of the picture is that the left path and the right path leave the
% same service and arrive at the same router, and only one of them involves the
% router talking to a subgraph at all. Everything below the router is what the
% two arrangements disagree about; everything above it is identical, which is
% why a client cannot tell them apart.
\begin{tikzpicture}[
    font=\small,
    box/.style={draw, rounded corners=2pt, minimum width=20mm,
                minimum height=7mm, align=center, font=\scriptsize},
    flow/.style={-{Stealth[length=2mm]}},
    held/.style={-{Stealth[length=2mm]}, thick},
    note/.style={font=\scriptsize, gray, align=center, inner sep=1pt}
  ]

  \node[box] (client)   at (0.4,0)    {client};
  \node[box] (router)   at (5.0,0)    {router};
  \node[box, minimum width=16mm] (accounts) at (9.8,0) {\code{accounts}};

  \node[box, fill=black!8] (reviews) at (1.9,-3.4) {\code{reviews}};
  \node[box, fill=black!8] (nats)    at (8.1,-3.4) {NATS};

  % -- the half a client can see ---------------------------------------------
  \draw[flow] (client) -- (router);
  \node[note, anchor=south, align=center] at (2.7,0.15)
    {two subscriptions,\\one write};

  \draw[flow] (router) -- (accounts);
  \node[note, anchor=south, align=center] at (7.4,0.15)
    {\code{\_entities},\\once per event};

  % -- left: the subscription the subgraph holds ------------------------------
  \draw[held] (router) -- (reviews);
  \node[note, anchor=east, align=right] at (2.55,-1.5)
    {one WebSocket,\\open as long as\\the client is};

  % -- right: the subscription the broker holds -------------------------------
  \draw[flow] (reviews) -- (nats);
  \node[note, anchor=south, align=center, fill=white, inner sep=2pt] at (5.0,-3.42)
    {\code{\_\_typename} and \code{id},\\published on a subject};

  \draw[flow] (nats) -- (router);
  \node[note, anchor=west, align=left] at (7.45,-1.5)
    {one subject,\\no service\\in between};

  \node[note, anchor=north, align=center] at (1.9,-3.85)
    {the write lands here,\\and announces itself twice};

  \node[note, anchor=north, align=center] at (8.1,-3.85)
    {\code{schema/streams.graphql}\\has no process behind it};

\end{tikzpicture}

  \caption{One write, announced twice, and the halves that differ are all below
  the router. On the left the router holds a connection to a service for as
  long as a client is listening; on the right it holds a subscription to a
  broker and the service that published is not involved in delivery at all. A
  client sees the same shape either way, which is what makes this a deployment
  decision rather than an API one.}
  \label{fig:ch14-two-arrangements}
\end{figure}

\subsection*{One word, and what its absence costs}

\index{event-driven federated subscriptions!the missing type name}%
Publish the same message without \code{\_\_typename}, so the body is only the
identifier, and the client gets this:

\begin{minted}[fontsize=\footnotesize,breaklines]{text}
{"errors":[{"message":"Cannot return null for non-nullable field
 'Subscription.reviewPublished.rating'."}],"data":null}
\end{minted}

The router received the message. Its debug log shows it arriving on the right
subject with the right body. It then makes no entity fetch and says nothing
about why, and the error names a field the publisher never mentioned and has no
opinion about. Nothing in the router, the subgraph, the schema or the composer
connects that error to the two missing tokens that caused it.

So the publisher carries a comment saying so, because nothing else would. That
is the same rule chapter~\ref{ch:hard-modeling-problems} reached about shared
scalars, applied one level down: the further a contract is from anything a tool
can check, the more of it has to be written in prose that people read.

\subsection*{What Mosaic had to build, which is one class}

\index{Mosaic.Reviews@\code{Mosaic.Reviews}!the publisher}%
Reviews gains a publisher and nothing else. The interesting line is the
subject:

\begin{minted}[fontsize=\footnotesize,breaklines]{csharp}
var subject = SubjectPrefix + accessor.Serializer.Format(nameof(Product), productId);
\end{minted}

\index{global object identifier!encoded rather than decoded}%
because the template interpolates the argument as the client wrote it, and a
client writes a global object identifier rather than a \code{Guid}. So this is
the one file in the repository that has to \emph{encode} a product identifier,
where the four copies of \code{ProductKey} that
chapter~\ref{ch:the-first-cut-extracting-catalog} argued about all decode one.
It reaches the serialiser the same way they do, through
\code{IResolverContext.Schema.Services}.

The mutation now announces the same fact twice. Both calls sit at the end of
the resolver, with a paragraph of comment between them, and stripped of that
they are:

\begin{minted}[fontsize=\footnotesize,breaklines]{csharp}
await sender.SendAsync(
    ReviewTopics.ReviewAdded(productId),
    review,
    cancellationToken);

await stream.PublishAsync(review, productId, context, cancellationToken);
\end{minted}

I left both, not because a real service would, but because the two
arrangements are this chapter's argument and a reader needs to watch them
answer the same write.

\index{contracts!a string two files have to agree about}%
One last thing to be nervous about. The publisher builds
\code{mosaic.reviewAdded.} and a template in a schema file says
\code{mosaic.reviewAdded.\{\{ args.productId \}\}}, and no compiler, composer or
type checker relates those two strings. It is a contract between a C\# constant
and a directive argument, and the only thing holding it together is that
\code{scripts/subscription-run.mjs} publishes through the mutation and asserts
that a frame comes out the other end. Chapter~\ref{ch:hard-modeling-problems}
said every shared type in a federated graph needs an owner and a written
definition outside both schemas. A shared subject name needs the same, and it
has less structure than a scalar.

\medskip
Two subscriptions then, built opposite ways, and a client cannot tell which is
which. What a client can see is a directive the graph advertises and does not
honour, which is the last thing in this chapter.

\index{event-driven federated subscriptions|)}%
